% Section 16.8, from lectures/L16/README.md.

\section{Summary}
\label{c16:sec:review}

\needspace{6\baselineskip}
\subsection*{What you should be able to do now}
\begin{itemize}
  \item Draw the CAN frame's state diagram and say which field every state corresponds to.
  \item Explain why every field needs a \code{LOAD} state before its shift state.
  \item Implement the transmit path: load the shift register, shift out, and move on at the right
    pulse.
  \item Drive the open-drain bus correctly: \code{bus_en} controls, \code{tx_bus} decides what.
  \item Gate the CRC engine correctly, and say why the wrong gating gives a CRC that looks plausible
    and is wrong.
\end{itemize}

\needspace{6\baselineskip}
\subsection*{Questions to test yourself with}
\begin{enumerate}
  \item Why does every field need a \code{LOAD} state of its own?
  \item A transmitting node releases the bus during the ACK slot. Why, and what does it expect to see
    there?
  \item \code{bus_en} and \code{tx_bus} are two separate signals. What would have been lost with just
    one?
  \item Why is \code{crc15} gated on \code{bit_valid and not stuff} and not on \code{bit_done}?
  \item Which fields does the CRC cover, and where exactly does the coverage end?
  \item Why is the state machine a Moore machine here?
\end{enumerate}

\needspace{6\baselineskip}
\subsection*{Target}
The transmit path should be in place when \chapref{c17:ch} begins, since the receive path is built
into the same machine. If you do not get all the way: finish the transmission through the CRC field,
and leave ACK and EOF to the next session.

\needspace{6\baselineskip}
\subsection*{Looking ahead}
\Chapref{c17:ch} completes \code{can_controller} with reception and arbitration, and runs the system
testbench all the way through.
